% Styles for the figures of "Attention as Matrix Algebra". Load with % Styles for the figures of "Attention as Matrix Algebra". Load with % Styles for the figures of "Attention as Matrix Algebra". Load with % Styles for the figures of "Attention as Matrix Algebra". Load with \input{attention-style}
% after \documentclass[tikz,border=2pt]{standalone}. The colours and line styles follow the
% diagram style of the Uncertainty Quantification project, copied here so the folder compiles
% on its own. No date or ID in the PDF, so a rebuild of an unchanged figure changes no file.
\pdfinfoomitdate=1
\pdftrailerid{}
\usepackage{helvet}
\renewcommand{\familydefault}{\sfdefault}
\usepackage{amsmath}
\usepackage{amssymb}
\usetikzlibrary{arrows.meta,calc,decorations.pathreplacing,fit,positioning,matrix}

% Colours, each with one role in every figure.
\definecolor{ink}{HTML}{172B3A}          % text, symbols and connectors
\definecolor{muted}{HTML}{62748E}        % short words under symbols
\definecolor{neutral}{HTML}{F1F5F9}      \definecolor{neutralline}{HTML}{62748E}
\definecolor{query}{HTML}{DFF2FE}        \definecolor{queryline}{HTML}{00A6F4}     % queries: sky blue
\definecolor{key}{HTML}{FEF3C6}          \definecolor{keyline}{HTML}{FFB900}       % keys: amber
\definecolor{value}{HTML}{D0FAE5}        \definecolor{valueline}{HTML}{00D492}     % values: green
\definecolor{pairs}{HTML}{FFE4E8}        \definecolor{pairsline}{HTML}{FF637E}     % n x n: pink
\definecolor{state}{HTML}{EDE9FE}        \definecolor{stateline}{HTML}{7E6CFF}     % fixed size: lavender
\colorlet{output}{neutral}               \colorlet{outputline}{neutralline}        % outputs: slate
\definecolor{crosstalk}{HTML}{FFEDD4}    \definecolor{crosstalkline}{HTML}{FF8904} % interference: orange
% A pseudo-value of the delta rule is a corrected value: a deeper green.
\definecolor{pseudo}{HTML}{8EECC4}       \colorlet{pseudoline}{valueline!75!ink}

% One text size: 7 pt, with Computer Modern math at the same size.
\newcommand{\seven}{\fontsize{7}{8.5}\selectfont}

% One cell of every drawn matrix, in cm.
\def\cell{0.34}

% \grid{name}{x}{y}{rows}{cols}{role}{highlights}: a matrix with its top-left corner at (x, y).
% Inside the highlights one unit is one cell, and row i spans y = -(i - 1) to y = -i.
\newcommand{\grid}[7]{%
  \begin{scope}[shift={(#2,#3)}, x=\cell cm, y=\cell cm]
    \fill[#6] (0,0) rectangle (#5,-#4);
    #7
    \draw[#6line, line width=0.3pt, step=1] (0,0) grid (#5,-#4);
    \draw[#6line, line width=0.8pt] (0,0) rectangle (#5,-#4);
    \node[inner sep=0pt, fit={(0,0) (#5,-#4)}] (#1) {};
  \end{scope}}
% Inside \grid: \row{colour}{i}{cols}, \col{colour}{j}{rows} and \entry{colour}{i}{j}.
\newcommand{\row}[3]{\fill[#1] (0,{-(#2)+1}) rectangle (#3,-#2);}
\newcommand{\col}[3]{\fill[#1] ({(#2)-1},0) rectangle (#2,-#3);}
\newcommand{\entry}[3]{\fill[#1] ({(#3)-1},{-(#2)+1}) rectangle (#3,-#2);}

\tikzset{
  every node/.append style={text=ink, font=\seven},
  tag/.style={line width=0.5pt, rounded corners=2pt, inner xsep=3pt, inner ysep=1.5pt,
    minimum height=4mm},
  title/.style={font=\seven\bfseries, inner sep=1pt},
  sym/.style={inner sep=1pt},
  word/.style={text=muted, align=center, inner sep=1pt},
  flow/.style={-{Stealth[length=2.4mm, width=2mm]}, line width=1.2pt, draw=ink, shorten >=1.5pt},
  branch/.style={-{Stealth[length=2mm, width=1.7mm]}, line width=0.9pt, draw=ink, shorten >=1pt},
  route/.style={line width=1.2pt, draw=ink},
  frame/.style={draw=neutralline, line width=0.6pt, rounded corners=5pt},
  matname/.style={sym, align=center},
  note/.style={word, align=center},
  subtitle/.style={word, font=\seven\itshape},
  keybox/.style={draw=neutralline, line width=0.5pt, dash pattern=on 1pt off 1.5pt,
    rounded corners=3pt, inner sep=4pt},
  chip/.style={tag, minimum width=5mm},
  rule/.style={draw=neutralline, line width=0.6pt, dash pattern=on 3pt off 2pt},
  % A small matrix written out with its numbers, one role colour: \matrix[entries=key] {1 \\ 0 \\};
  entries/.style={matrix of math nodes, column sep=-0.4pt, row sep=-0.4pt,
    nodes={draw=#1line, fill=#1, line width=0.4pt, minimum size=5mm, inner sep=0pt, anchor=center}},
}

% after \documentclass[tikz,border=2pt]{standalone}. The colours and line styles follow the
% diagram style of the Uncertainty Quantification project, copied here so the folder compiles
% on its own. No date or ID in the PDF, so a rebuild of an unchanged figure changes no file.
\pdfinfoomitdate=1
\pdftrailerid{}
\usepackage{helvet}
\renewcommand{\familydefault}{\sfdefault}
\usepackage{amsmath}
\usepackage{amssymb}
\usetikzlibrary{arrows.meta,calc,decorations.pathreplacing,fit,positioning,matrix}

% Colours, each with one role in every figure.
\definecolor{ink}{HTML}{172B3A}          % text, symbols and connectors
\definecolor{muted}{HTML}{62748E}        % short words under symbols
\definecolor{neutral}{HTML}{F1F5F9}      \definecolor{neutralline}{HTML}{62748E}
\definecolor{query}{HTML}{DFF2FE}        \definecolor{queryline}{HTML}{00A6F4}     % queries: sky blue
\definecolor{key}{HTML}{FEF3C6}          \definecolor{keyline}{HTML}{FFB900}       % keys: amber
\definecolor{value}{HTML}{D0FAE5}        \definecolor{valueline}{HTML}{00D492}     % values: green
\definecolor{pairs}{HTML}{FFE4E8}        \definecolor{pairsline}{HTML}{FF637E}     % n x n: pink
\definecolor{state}{HTML}{EDE9FE}        \definecolor{stateline}{HTML}{7E6CFF}     % fixed size: lavender
\colorlet{output}{neutral}               \colorlet{outputline}{neutralline}        % outputs: slate
\definecolor{crosstalk}{HTML}{FFEDD4}    \definecolor{crosstalkline}{HTML}{FF8904} % interference: orange
% A pseudo-value of the delta rule is a corrected value: a deeper green.
\definecolor{pseudo}{HTML}{8EECC4}       \colorlet{pseudoline}{valueline!75!ink}

% One text size: 7 pt, with Computer Modern math at the same size.
\newcommand{\seven}{\fontsize{7}{8.5}\selectfont}

% One cell of every drawn matrix, in cm.
\def\cell{0.34}

% \grid{name}{x}{y}{rows}{cols}{role}{highlights}: a matrix with its top-left corner at (x, y).
% Inside the highlights one unit is one cell, and row i spans y = -(i - 1) to y = -i.
\newcommand{\grid}[7]{%
  \begin{scope}[shift={(#2,#3)}, x=\cell cm, y=\cell cm]
    \fill[#6] (0,0) rectangle (#5,-#4);
    #7
    \draw[#6line, line width=0.3pt, step=1] (0,0) grid (#5,-#4);
    \draw[#6line, line width=0.8pt] (0,0) rectangle (#5,-#4);
    \node[inner sep=0pt, fit={(0,0) (#5,-#4)}] (#1) {};
  \end{scope}}
% Inside \grid: \row{colour}{i}{cols}, \col{colour}{j}{rows} and \entry{colour}{i}{j}.
\newcommand{\row}[3]{\fill[#1] (0,{-(#2)+1}) rectangle (#3,-#2);}
\newcommand{\col}[3]{\fill[#1] ({(#2)-1},0) rectangle (#2,-#3);}
\newcommand{\entry}[3]{\fill[#1] ({(#3)-1},{-(#2)+1}) rectangle (#3,-#2);}

\tikzset{
  every node/.append style={text=ink, font=\seven},
  tag/.style={line width=0.5pt, rounded corners=2pt, inner xsep=3pt, inner ysep=1.5pt,
    minimum height=4mm},
  title/.style={font=\seven\bfseries, inner sep=1pt},
  sym/.style={inner sep=1pt},
  word/.style={text=muted, align=center, inner sep=1pt},
  flow/.style={-{Stealth[length=2.4mm, width=2mm]}, line width=1.2pt, draw=ink, shorten >=1.5pt},
  branch/.style={-{Stealth[length=2mm, width=1.7mm]}, line width=0.9pt, draw=ink, shorten >=1pt},
  route/.style={line width=1.2pt, draw=ink},
  frame/.style={draw=neutralline, line width=0.6pt, rounded corners=5pt},
  matname/.style={sym, align=center},
  note/.style={word, align=center},
  subtitle/.style={word, font=\seven\itshape},
  keybox/.style={draw=neutralline, line width=0.5pt, dash pattern=on 1pt off 1.5pt,
    rounded corners=3pt, inner sep=4pt},
  chip/.style={tag, minimum width=5mm},
  rule/.style={draw=neutralline, line width=0.6pt, dash pattern=on 3pt off 2pt},
  % A small matrix written out with its numbers, one role colour: \matrix[entries=key] {1 \\ 0 \\};
  entries/.style={matrix of math nodes, column sep=-0.4pt, row sep=-0.4pt,
    nodes={draw=#1line, fill=#1, line width=0.4pt, minimum size=5mm, inner sep=0pt, anchor=center}},
}

% after \documentclass[tikz,border=2pt]{standalone}. The colours and line styles follow the
% diagram style of the Uncertainty Quantification project, copied here so the folder compiles
% on its own. No date or ID in the PDF, so a rebuild of an unchanged figure changes no file.
\pdfinfoomitdate=1
\pdftrailerid{}
\usepackage{helvet}
\renewcommand{\familydefault}{\sfdefault}
\usepackage{amsmath}
\usepackage{amssymb}
\usetikzlibrary{arrows.meta,calc,decorations.pathreplacing,fit,positioning,matrix}

% Colours, each with one role in every figure.
\definecolor{ink}{HTML}{172B3A}          % text, symbols and connectors
\definecolor{muted}{HTML}{62748E}        % short words under symbols
\definecolor{neutral}{HTML}{F1F5F9}      \definecolor{neutralline}{HTML}{62748E}
\definecolor{query}{HTML}{DFF2FE}        \definecolor{queryline}{HTML}{00A6F4}     % queries: sky blue
\definecolor{key}{HTML}{FEF3C6}          \definecolor{keyline}{HTML}{FFB900}       % keys: amber
\definecolor{value}{HTML}{D0FAE5}        \definecolor{valueline}{HTML}{00D492}     % values: green
\definecolor{pairs}{HTML}{FFE4E8}        \definecolor{pairsline}{HTML}{FF637E}     % n x n: pink
\definecolor{state}{HTML}{EDE9FE}        \definecolor{stateline}{HTML}{7E6CFF}     % fixed size: lavender
\colorlet{output}{neutral}               \colorlet{outputline}{neutralline}        % outputs: slate
\definecolor{crosstalk}{HTML}{FFEDD4}    \definecolor{crosstalkline}{HTML}{FF8904} % interference: orange
% A pseudo-value of the delta rule is a corrected value: a deeper green.
\definecolor{pseudo}{HTML}{8EECC4}       \colorlet{pseudoline}{valueline!75!ink}

% One text size: 7 pt, with Computer Modern math at the same size.
\newcommand{\seven}{\fontsize{7}{8.5}\selectfont}

% One cell of every drawn matrix, in cm.
\def\cell{0.34}

% \grid{name}{x}{y}{rows}{cols}{role}{highlights}: a matrix with its top-left corner at (x, y).
% Inside the highlights one unit is one cell, and row i spans y = -(i - 1) to y = -i.
\newcommand{\grid}[7]{%
  \begin{scope}[shift={(#2,#3)}, x=\cell cm, y=\cell cm]
    \fill[#6] (0,0) rectangle (#5,-#4);
    #7
    \draw[#6line, line width=0.3pt, step=1] (0,0) grid (#5,-#4);
    \draw[#6line, line width=0.8pt] (0,0) rectangle (#5,-#4);
    \node[inner sep=0pt, fit={(0,0) (#5,-#4)}] (#1) {};
  \end{scope}}
% Inside \grid: \row{colour}{i}{cols}, \col{colour}{j}{rows} and \entry{colour}{i}{j}.
\newcommand{\row}[3]{\fill[#1] (0,{-(#2)+1}) rectangle (#3,-#2);}
\newcommand{\col}[3]{\fill[#1] ({(#2)-1},0) rectangle (#2,-#3);}
\newcommand{\entry}[3]{\fill[#1] ({(#3)-1},{-(#2)+1}) rectangle (#3,-#2);}

\tikzset{
  every node/.append style={text=ink, font=\seven},
  tag/.style={line width=0.5pt, rounded corners=2pt, inner xsep=3pt, inner ysep=1.5pt,
    minimum height=4mm},
  title/.style={font=\seven\bfseries, inner sep=1pt},
  sym/.style={inner sep=1pt},
  word/.style={text=muted, align=center, inner sep=1pt},
  flow/.style={-{Stealth[length=2.4mm, width=2mm]}, line width=1.2pt, draw=ink, shorten >=1.5pt},
  branch/.style={-{Stealth[length=2mm, width=1.7mm]}, line width=0.9pt, draw=ink, shorten >=1pt},
  route/.style={line width=1.2pt, draw=ink},
  frame/.style={draw=neutralline, line width=0.6pt, rounded corners=5pt},
  matname/.style={sym, align=center},
  note/.style={word, align=center},
  subtitle/.style={word, font=\seven\itshape},
  keybox/.style={draw=neutralline, line width=0.5pt, dash pattern=on 1pt off 1.5pt,
    rounded corners=3pt, inner sep=4pt},
  chip/.style={tag, minimum width=5mm},
  rule/.style={draw=neutralline, line width=0.6pt, dash pattern=on 3pt off 2pt},
  % A small matrix written out with its numbers, one role colour: \matrix[entries=key] {1 \\ 0 \\};
  entries/.style={matrix of math nodes, column sep=-0.4pt, row sep=-0.4pt,
    nodes={draw=#1line, fill=#1, line width=0.4pt, minimum size=5mm, inner sep=0pt, anchor=center}},
}

% after \documentclass[tikz,border=2pt]{standalone}. The colours and line styles follow the
% diagram style of the Uncertainty Quantification project, copied here so the folder compiles
% on its own. No date or ID in the PDF, so a rebuild of an unchanged figure changes no file.
\pdfinfoomitdate=1
\pdftrailerid{}
\usepackage{helvet}
\renewcommand{\familydefault}{\sfdefault}
\usepackage{amsmath}
\usepackage{amssymb}
\usetikzlibrary{arrows.meta,calc,decorations.pathreplacing,fit,positioning,matrix}

% Colours, each with one role in every figure.
\definecolor{ink}{HTML}{172B3A}          % text, symbols and connectors
\definecolor{muted}{HTML}{62748E}        % short words under symbols
\definecolor{neutral}{HTML}{F1F5F9}      \definecolor{neutralline}{HTML}{62748E}
\definecolor{query}{HTML}{DFF2FE}        \definecolor{queryline}{HTML}{00A6F4}     % queries: sky blue
\definecolor{key}{HTML}{FEF3C6}          \definecolor{keyline}{HTML}{FFB900}       % keys: amber
\definecolor{value}{HTML}{D0FAE5}        \definecolor{valueline}{HTML}{00D492}     % values: green
\definecolor{pairs}{HTML}{FFE4E8}        \definecolor{pairsline}{HTML}{FF637E}     % n x n: pink
\definecolor{state}{HTML}{EDE9FE}        \definecolor{stateline}{HTML}{7E6CFF}     % fixed size: lavender
\colorlet{output}{neutral}               \colorlet{outputline}{neutralline}        % outputs: slate
\definecolor{crosstalk}{HTML}{FFEDD4}    \definecolor{crosstalkline}{HTML}{FF8904} % interference: orange
% A pseudo-value of the delta rule is a corrected value: a deeper green.
\definecolor{pseudo}{HTML}{8EECC4}       \colorlet{pseudoline}{valueline!75!ink}

% One text size: 7 pt, with Computer Modern math at the same size.
\newcommand{\seven}{\fontsize{7}{8.5}\selectfont}

% One cell of every drawn matrix, in cm.
\def\cell{0.34}

% \grid{name}{x}{y}{rows}{cols}{role}{highlights}: a matrix with its top-left corner at (x, y).
% Inside the highlights one unit is one cell, and row i spans y = -(i - 1) to y = -i.
\newcommand{\grid}[7]{%
  \begin{scope}[shift={(#2,#3)}, x=\cell cm, y=\cell cm]
    \fill[#6] (0,0) rectangle (#5,-#4);
    #7
    \draw[#6line, line width=0.3pt, step=1] (0,0) grid (#5,-#4);
    \draw[#6line, line width=0.8pt] (0,0) rectangle (#5,-#4);
    \node[inner sep=0pt, fit={(0,0) (#5,-#4)}] (#1) {};
  \end{scope}}
% Inside \grid: \row{colour}{i}{cols}, \col{colour}{j}{rows} and \entry{colour}{i}{j}.
\newcommand{\row}[3]{\fill[#1] (0,{-(#2)+1}) rectangle (#3,-#2);}
\newcommand{\col}[3]{\fill[#1] ({(#2)-1},0) rectangle (#2,-#3);}
\newcommand{\entry}[3]{\fill[#1] ({(#3)-1},{-(#2)+1}) rectangle (#3,-#2);}

\tikzset{
  every node/.append style={text=ink, font=\seven},
  tag/.style={line width=0.5pt, rounded corners=2pt, inner xsep=3pt, inner ysep=1.5pt,
    minimum height=4mm},
  title/.style={font=\seven\bfseries, inner sep=1pt},
  sym/.style={inner sep=1pt},
  word/.style={text=muted, align=center, inner sep=1pt},
  flow/.style={-{Stealth[length=2.4mm, width=2mm]}, line width=1.2pt, draw=ink, shorten >=1.5pt},
  branch/.style={-{Stealth[length=2mm, width=1.7mm]}, line width=0.9pt, draw=ink, shorten >=1pt},
  route/.style={line width=1.2pt, draw=ink},
  frame/.style={draw=neutralline, line width=0.6pt, rounded corners=5pt},
  matname/.style={sym, align=center},
  note/.style={word, align=center},
  subtitle/.style={word, font=\seven\itshape},
  keybox/.style={draw=neutralline, line width=0.5pt, dash pattern=on 1pt off 1.5pt,
    rounded corners=3pt, inner sep=4pt},
  chip/.style={tag, minimum width=5mm},
  rule/.style={draw=neutralline, line width=0.6pt, dash pattern=on 3pt off 2pt},
  % A small matrix written out with its numbers, one role colour: \matrix[entries=key] {1 \\ 0 \\};
  entries/.style={matrix of math nodes, column sep=-0.4pt, row sep=-0.4pt,
    nodes={draw=#1line, fill=#1, line width=0.4pt, minimum size=5mm, inner sep=0pt, anchor=center}},
}
